\documentclass[11pt]{article}
\usepackage[margin=1in]{geometry}
\usepackage{amsmath,amssymb,amsthm,hyperref}
\newcommand{\Q}{\mathbb Q}
\newcommand{\CH}{\mathrm{CH}}
\newcommand{\Hdg}{\mathrm{Hdg}}
\newcommand{\KU}{\mathrm{KU}}
\newcommand{\Perf}{\mathrm{Perf}}
\newcommand{\cl}{\mathrm{cl}}
\newcommand{\ch}{\mathrm{ch}}
\newcommand{\td}{\mathrm{td}}
\newcommand{\real}{\mathrm{real}}
\newcommand{\im}{\mathrm{im}}
\newcommand{\id}{\mathrm{id}}
\newtheorem{theorem}{Theorem}
\newtheorem{lemma}[theorem]{Lemma}
\newtheorem{proposition}[theorem]{Proposition}
\newtheorem{problem}[theorem]{Unresolved geometric problem}
\title{GST native residual Euler kernels and complete-alteration transfer}
\author{Mathbox: mathematical construction and source audit}
\date{20 October 2026}
\begin{document}
\maketitle
\begin{abstract}
Starting from the established GST normal-free middle-world Euler law,
we construct an actual rational algebraic perfect-kernel class that
removes any finite native primitive frame. Its Euler pairing against
a doubled prescribed topological class is exactly minus the positive
squared residual. The construction works for every original variety
in the stated middle-primitive range and survives every smooth
complete alteration, with the exact degree factor. We derive a
native-source-to-original-cycle formula for the remaining linear
detector. The universal existence of that linear native source is
not proved. In particular, a native kernel detecting a doubled
topological input is not a native representative of the input.
\end{abstract}

\section{Scope, pinned input and conventions}
The source revision is
\texttt{c2479cac309ce2379c47e77b912f93183923e9b2}
of \texttt{kyo-oo/HC-PROOF}. The established input is the
normalization theorem in
\texttt{GST\_middle\_Euler\_arithmetic\_normalization\_2026-10-20.tex}.
No new Lean theorem or kernel compilation is claimed in this note.

Let $X$ be a connected smooth projective complex variety of dimension
$d$. Embed $X$ in $\mathbf P^{2M}$, increasing $M$ if necessary so
$M\ge d+1$. Put
\[
 Y=\operatorname{Bl}_X\mathbf P^{2M},\qquad
 j:E=\mathbf P_{\mathrm{lines}}(N)\hookrightarrow Y,\qquad
 q:E\longrightarrow X,\qquad c=2M-d .
\]
Our convention is $\mathcal O_E(-1)\subset q^*N$,
$\xi=c_1(\mathcal O_E(1))$, and $s(N)=c(N)^{-1}$.
Set
\[
 \mathcal H_X=\bigoplus_{a=0}^d\Hdg^a(X)_\Q,\qquad
 P_X(x,y)=\int_X[xy]_d,\qquad \epsilon=(-1)^{M+1}.
\]
The previously constructed finite native characteristic-class word
$R=\sqrt{c(N)}$ satisfies $R^2s(N)=1$.
The rational topological class $A_x=\widehat V_x\in\KU^0(Y)_\Q$
is linear in $x$ and has
\[
 \ch(A_x)=I(Rx),\qquad
 I(x)=\sum_{a=0}^d
 j_*(\xi^{M-a-1}q^*x_a),\qquad
 \chi_Y(A_x,A_y)=\epsilon P_X(x,y).                 \tag{1}
\]
The entire character is concentrated in codimension $M$.
Consequently $A_x^\vee=(-1)^M A_x$ in rational topological $K$-theory.

For the positive residual construction take $d=2p$, $p\ge1$, and
an ample rational class $h$. Define
\[
 V=\ker(h:H^{2p}(X,\Q)\to H^{2p+2}(X,\Q))
       \cap H^{p,p}(X),\qquad
 Q(x,y)=(-1)^p\int_Xxy .
\]
The established primitive Hodge--Riemann form makes $Q$ positive
definite on $V$. On this sector (1) becomes
\[
 \chi_Y(A_x,A_y)=\sigma Q(x,y),\qquad
 \sigma=(-1)^{M+p+1},\qquad A_{Jx}=(-1)^p A_x,
 \quad Jx=(-1)^p x.                                \tag{2}
\]
Positivity is used only on $V$, or on its injective alteration image.

\section{Actual native lifts of a known primitive frame}
Choose actual rational cycles $z_1,\ldots,z_t\in\CH^p(X)_\Q$
whose classes $v_i=\cl(z_i)\in V$ are linearly independent.
This is any already constructed native frame, not a basis assumed
to cover all of $V$. Define actual middle Chow classes
\[
 \Gamma_i=I^{\CH}(R^{\CH}z_i)\in\CH^M(Y)_\Q .
\]
The rational algebraic Chern-character isomorphism for smooth $Y$
gives classes $C_i\in K_0^{\mathrm{alg}}(Y)_\Q$ with
$\ch_{\mathrm{alg}}(C_i)=\Gamma_i$ and all other character degrees zero.
Hence
\[
 \real C_i=A_{v_i},\qquad C_i^\vee=(-1)^M C_i.
\]
These are genuine rational algebraic $K$-classes. Clearing finite
denominators realizes them by bounded perfect complexes using finite
direct sums and shifts; this does not assert that a rational class
is literally one integral object without clearing denominators.

Let
\[
 B_{ij}=Q(v_i,v_j),\qquad
 G_{ij}=\chi_Y(C_i,C_j)=\sigma B_{ij}.                \tag{3}
\]
Independence and positivity imply $\det B>0$, so $G$ is invertible.
For any prescribed $\alpha\in V$, define
\[
 q_i=Q(v_i,\alpha),\qquad
 \alpha_0=\sum_jv_j(B^{-1}q)_j,\qquad
 r=\alpha-\alpha_0.                                 \tag{4}
\]
These are rational classes; $Q(v_i,r)=0$.
The actual cycle $\sum_j(B^{-1}q)_jz_j$ represents $\alpha_0$.
No actual cycle representing $r$ is introduced.

\section{Constructing the native residual kernel}
Write $p_1,p_2:Y\times Y\to Y$ for the projections.
For a rational algebraic perfect-kernel class $F$, put
\[
 T_F(U)=(p_2)_!\bigl(p_1^*U\cdot F\bigr).
\]
This is defined on algebraic $K_0$ and on topological $K$-theory;
realization commutes with it. Tensor products and proper $K$-pushforward
are used. The kernel $[\mathcal O_{\Delta_Y}]$ is the identity.

\begin{theorem}[Actual finite-frame and residual kernels]
The following are actual rational algebraic perfect-kernel classes:
\[
 \mathscr P=\sum_{i,j}(G^{-1})_{ji}\,C_i^\vee\boxtimes C_j,
 \qquad
 \boxed{\mathscr Q=[\mathcal O_{\Delta_Y}]-\mathscr P.} \tag{5}
\]
For convolution of kernels,
\[
 \mathscr P\star\mathscr P=\mathscr P,\qquad
 \mathscr Q\star\mathscr Q=\mathscr Q,\qquad
 \mathscr P\star\mathscr Q=\mathscr Q\star\mathscr P=0. \tag{6}
\]
On the prescribed middle classes,
\[
 T_{\mathscr P}(A_\alpha)=A_{\alpha_0},\qquad
 \boxed{T_{\mathscr Q}(A_\alpha)=A_r.}                \tag{7}
\]
Moreover $T_{\mathscr Q}(C_i)=0$ in algebraic $K_0(Y)_\Q$.
\end{theorem}
\begin{proof}
External product and proper pushforward give
\[
 T_{\mathscr P}(U)
 =\sum_{i,j}(G^{-1})_{ji}\chi_Y(C_i,U)C_j.             \tag{8}
\]
This identity holds in the algebraic theory for algebraic $U$,
and after realization for topological $U$.
The convolution of two decomposable kernels contracts the middle
factors by $\chi_Y(C_k,C_j)=G_{kj}$. Thus the coefficient matrix
for the square is $G^{-1}GG^{-1}=G^{-1}$, proving the first
identity in (6) at the level of kernel classes, not only on
their cohomological actions. Subtracting from the identity kernel
gives the remaining identities.

For $U=A_\alpha$, the period column in (8) is $\sigma q$.
Since $G^{-1}=\sigma B^{-1}$, (8) is precisely the first
identity of (7). Linearity gives the second.
For $U=C_k$, (8) gives
$\sum_j(G^{-1}G)_{jk}C_j=C_k$, proving the final assertion.
\end{proof}
When $t=0$ use $\mathscr P=0$, $\mathscr Q=\mathcal O_\Delta$,
$\alpha_0=0$. No inverse of a nonempty singular matrix is used.

\section{A universal native quadratic detector}
The next identity is a construction of a native kernel with a
nonzero index. Its test object is topological and doubled.

\begin{lemma}[Diagonal and frame contractions]
For pure-middle rational topological classes $A,B$,
\[
 \chi_{Y\times Y}(A\boxtimes B,\mathcal O_\Delta)
       =(-1)^M\chi_Y(A,B).                           \tag{9}
\]
If $u(A)_i=\chi_Y(C_i,A)$ and likewise for $B$, then
\[
 \chi_{Y\times Y}(A\boxtimes B,\mathscr Q)
 =(-1)^M\left(\chi_Y(A,B)
              -u(B)^TG^{-1}u(A)\right).             \tag{10}
\]
\end{lemma}
\begin{proof}
Grothendieck--Riemann--Roch for the diagonal gives the left side
of (9) as $\int_Y\ch(A^\vee)\ch(B^\vee)\td(Y)$.
Both characters have degree $M$, so only $\td_0=1$ contributes.
The two duality signs cancel, leaving
$\int_Y\ch(A)\ch(B)=(-1)^M\chi_Y(A,B)$.
For (10), Euler pairing factors over external products.
Since $C_i^\vee=(-1)^MC_i$ and the pure-middle Euler pairing
is symmetric,
\[
 \chi_{Y\times Y}(A\boxtimes B,C_i^\vee\boxtimes C_j)
       =(-1)^M u(A)_iu(B)_j.
\]
Sum with the coefficients in (5) and subtract from (9).
\end{proof}

\begin{theorem}[Exact residual-energy law]
For every $\alpha\in V$ and every independently constructed
native frame as above, the actual kernel (5) satisfies
\[
 \boxed{
 -\chi_{Y\times Y}
       (A_\alpha\boxtimes A_{J\alpha},\mathscr Q)
       =Q(r,r)
       =Q(\alpha,\alpha)-q^TB^{-1}q.}                \tag{11}
\]
This is strictly positive if and only if $r\ne0$.
For the empty frame it holds for every nonzero $\alpha$ using
only the actual diagonal kernel.
\end{theorem}
\begin{proof}
In (10), $B=A_{J\alpha}=(-1)^p A_\alpha$,
$u(A_\alpha)=\sigma q$, and $G^{-1}=\sigma B^{-1}$,
where the matrix $B$ on the right is that of (3).
Using $\sigma=(-1)^p\epsilon$, the bracket in (10) becomes
\[
 \epsilon\bigl(Q(\alpha,\alpha)-q^TB^{-1}q\bigr)
       =\epsilon Q(r,r).
\]
Finally $(-1)^M\epsilon=-1$. Orthogonality in (4)
proves the squared-residual expression, and positive
definiteness gives strict positivity exactly when $r\ne0$.
\end{proof}
All quantities are rational; finite positive denominator clearing
makes the nonzero index integral. Such clearing changes no
algebraicity or nonvanishing assertion.

\section{Complete alterations preserve the construction with exact degree}
Let $f:X'\to X$ be any proper generically finite dominant morphism
of connected smooth projective varieties, with $\deg f=e>0$.
No flatness assumption is needed for operational Chow pullback
between the smooth varieties. Projection gives
\[
 f_*f^*=e\,\id,\qquad
 \int_{X'}f^*x\,f^*y=e\int_Xxy.                      \tag{12}
\]
In particular $f^*$ is injective. The form $(-1)^p\int_{X'}xy$
is positive definite on $f^*V$, although $f^*V$ need not be
primitive for any chosen polarization on $X'$.

Choose one $M$ large enough to embed both varieties in
$\mathbf P^{2M}$, and construct $Y'$ and the separately
normalized classes $A'_x$ using the normal bundle of $X'$.
Use the actual native frame $f^*z_i$. Its Gram matrices satisfy
\[
 B'=eB,\qquad G'=eG.
\]
The projection coefficients in (4) are unchanged because
$(B')^{-1}q'=(eB)^{-1}(eq)=B^{-1}q$.
Construct the actual kernel $\mathscr Q'$ by (5) on $Y'\times Y'$.

\begin{theorem}[Blowup--alteration residual-index law]
For every such complete alteration and every $\alpha\in V$,
\[
 \boxed{
 -\chi_{Y'\times Y'}
   (A'_{f^*\alpha}\boxtimes A'_{f^*J\alpha},\mathscr Q')
       =e\,Q(r,r).}                                 \tag{13}
\]
Thus neither the change of normal bundle nor alteration degree
destroys the residual detector. The result is exact for every
positive degree, not a finite-degree experiment.
\end{theorem}
\begin{proof}
Apply the proof of (11) on the positive subspace $f^*V$.
Equations (12) scale both the total norm and the Schur term
by $e$, leaving the projection coefficients unchanged.
The normal bundles do not enter after the independent
GST normalizations.
\end{proof}

\section{The exact remaining linear-source construction}
For any actual middle Chow class $\Gamma$ on $Y'$, the established
Chern--Segre returns are
\[
 T_a(\Gamma)=-q'_*\bigl((\xi')^{c'-1-a}j'^*\Gamma\bigr),\qquad
 \mathcal R_a(\Gamma)=
    \sum_{k=0}^{c'-1-a}c_k(N')T_{a+k}(\Gamma).
\]
Put
\[
 \mathfrak a'(\Gamma)
 =(R'^{\CH})^{-1}
    \sum_{b=0}^{2p}\mathcal R_{M-b}(\Gamma)
       \in\bigoplus_b\CH^b(X')_\Q.                  \tag{14}
\]
Every operator in (14) is actual Chow geometry. The sum is graded.
For an independently existing $F\in K_0^{\mathrm{alg}}(Y')_\Q$,
use its actual Adams filter
\[
 E=\Pi_MF,\qquad
 \Pi_M=\prod_{\substack{0\le k\le2M\\k\ne M}}
             \frac{\psi^2-2^k}{2^M-2^k},\qquad
 z'_E=[\mathfrak a'(\ch_M E)]_p.                    \tag{15}
\]
This processes a source already known to be algebraic.
It does not apply an algebraic Adams operation to an unknown
algebraic lift of a topological target.

\begin{proposition}[Actual original-$X$ source returned from an altered world]
For every $r\in V$, every such $f$, and every actual $E$ in (15),
the actual rational cycle
\[
 \boxed{Z_E=e^{-1}f_*z'_E\in\CH^p(X)_\Q}
\]
satisfies
\[
 \boxed{
 (-1)^p\int_Xr\,\cl(Z_E)
   =\frac{(-1)^{M+p+1}}e
       \chi_{Y'}(A'_{f^*r},\real E).}               \tag{16}
\]
\end{proposition}
\begin{proof}
Write $a'=\cl(\mathfrak a'(\ch_M E))$.
The canonical blowup splitting says
$\ch(\real E)=t\pi'^*H^M+I'(R'a')$.
The ambient term is orthogonal to $I'(R'f^*r)$ since
$H^M|_{X'}=0$ for $M>2p$.
Equation (1) then gives
\[
 \chi_{Y'}(A'_{f^*r},\real E)
      =\epsilon\int_{X'}f^*r\,a'_p .
\]
Projection and the definition of $Z_E$ yield (16).
A nonzero right side therefore supplies a genuine new original-$X$
cycle having nonzero residual period. Primitivity of that new
cycle is not asserted without a further actual purification.
\end{proof}

Formula (16) proves the blowup--complete-alteration return.
The unresolved input is an actual single-world algebraic $F$
for which the right side is nonzero for an as yet undetected
prescribed residual. Formula (13) supplies a doubled-world
kernel test, which has different input types.

\section{Why the quadratic construction is not yet a linear native source}
By (7), contracting the native kernel with the prescribed input gives
\[
 T_{\mathscr Q}(A_\alpha)=A_r\in\KU^0(Y)_\Q .
\]
It belongs to algebraic $K_0$ if and only if the required residual
has an algebraic lift under the established split normalization.
The contraction uses $p_1^*A_\alpha$, which is topological unless
its algebraicity has independently been proved.
Applying a genuine algebraic kernel does not change this input type.

On the other hand, for an actually algebraic input $D$,
$T_{\mathscr Q}(D)$ is algebraic; its residual pairing can be
nonzero only if $D$ already supplies a detected direction outside
the chosen frame. The kernel annihilates its known frame exactly.

The distinction survives tests on actual point objects.
For any closed point $y\in Y$,
\[
 \chi_Y(C_i,\mathcal O_y)=\operatorname{rank}(C_i)=0,\qquad
 T_{\mathscr Q}(\mathcal O_y)=\mathcal O_y,\qquad
 \Pi_M[\mathcal O_y]=0.
\]
The last equality follows since $\mathcal O_y$ has only
top-codimension $2M$ character and $M>0$.
Thus a nonzero kernel action on genuine objects, even all point
objects, does not supply a nonzero middle residual source.

\begin{proposition}[Exact test of the native partial-trace candidate]
For every actual algebraic input $D\in K_0^{\mathrm{alg}}(Y)_\Q$,
$F_D=T_{\mathscr Q}(D)$ is an actual algebraic source and
\[
 \boxed{\chi_Y(A_r,\real F_D)=\chi_Y(A_r,\real D).}
\]
The same equality of residual periods holds after applying
the actual middle Adams filter to both sources. In particular
\[
 \boxed{\chi_Y(A_r,\real(\Pi_MT_{\mathscr Q}(\mathcal O_Y)))=0.}
\]
The unfiltered unit partial trace does not improve the residual
period of $\mathcal O_Y$.
\end{proposition}
\begin{proof}
Equation (8) writes $T_{\mathscr P}(D)$ as a finite rational
combination of the $C_i$. Each has zero pairing with $A_r$
by (2) and $Q(r,v_i)=0$. Subtract from the identity action.
The filtered difference is the same finite combination because
$\Pi_MC_i=C_i$. Since $\Pi_M\mathcal O_Y=0$ for $M>0$,
the filtered unit period is zero.
For the unfiltered unit, the remaining period is exactly
\[
 \chi_Y(A_r,\mathcal O_Y)
   =(-1)^M\int_Y I(Rr)\,\td_M(Y).
\]
This is a single existing characteristic-class period, not the
positive quadratic form $Q(r,r)$.
\end{proof}
More generally, applying this residual kernel to a finite list
of native inputs preserves every one of their residual periods.
If those periods all vanish, the partial-trace construction
supplies no new detector. A fixed single source gives a linear
functional of $r$, whereas (11) is quadratic: substituting $-r$
reverses the former and preserves the latter. Thus they cannot
be identified as the same positive law on all nonzero residuals.

\section{The exact universal obligation and the existing Lean inputs}
For arbitrary dimension return to $\mathcal H_X$ in (1).
The following is an exact algebraic-source obligation:
\[
 \begin{gathered}
 \text{For every }0\ne x\in\mathcal H_X,\text{ construct }
 F\in K_0^{\mathrm{alg}}(Y)_\Q\text{ such that}\\
 \boxed{\chi_Y(A_x,\real(\Pi_MF))\ne0.}              \tag{17}
 \end{gathered}
\]
This asks for actual independent source geometry; choosing
$F=A_{Jx}$ in topological $K$-theory does not satisfy it.

\begin{proposition}[Exact status of the universal source obligation]
For a fixed original $X$, (17) is equivalent to all rational
Hodge classes on $X$ being algebraic. Allowing smooth complete
alterations and the return (14) does not weaken this implication.
\end{proposition}
\begin{proof}
Let $\mathcal A_X$ be the total image of canonical rational Chow
cycle classes. Actual returns (14) lie in $\mathcal A_X$.
By (1), the linear Euler tests in (17) are precisely, up to
the nonzero sign $\epsilon$, the functionals $P_X(x,a)$ for
actual $a\in\mathcal A_X$. The full rational Hodge pairing
$P_X$ is perfect by the established rational polarization.
Consequently all nonzero $x$ being separated is equivalent to
$\mathcal A_X^\perp=0$, hence to $\mathcal A_X=\mathcal H_X$
by finite-dimensional duality. No nondegeneracy assumption
on the restriction to $\mathcal A_X$ is used.

Conversely, if all Hodge classes are algebraic, choose an
algebraic packet $a$ with $P_X(x,a)\ne0$ and its actual native
normalized lift. This gives (17).
For alterations, projection gives
$\int_{X'}f^*x\,a'=\int_Xx\,f_*a'$ with the understood graded
complementary degrees. Every $f_*a'$ is still algebraic on $X$.
The identity alteration supplies the converse.
\end{proof}

The pinned Lean declarations agree with this precise boundary:
\begin{itemize}
\item \texttt{NativeArsenalRepresentation.hodge\_action} supplies
native matrix-unit actions. The global arsenal theorem takes
this representation as an input; it does not construct it on all $X$.
\item \texttt{FullGSTPlaneCompleteness} asks for an actual native
program for every Hodge basis direction.
\texttt{hodge\_of\_fullGSTPlane} consumes that hypothesis.
\item \texttt{CMGSTNativeSheetProgram.basisPeriodExpansion}
specifies that the actual transferred native-source periods span
the requested basis direction.
\texttt{CMGSTNativeBasisProgramSupply} asserts existence of
these geometric packets for every direction.
\item The intrinsic \texttt{gst\_plane\_completeness} theorem
constructs coordinate/causal branches unconditionally.
The same file expressly separates it from native geometric landing.
\end{itemize}
None of these compiled laws is contradicted by this note.
Their native representation or source-supply inputs are exactly
where a proof of (17) would have to enter.

\section{Exact computation and self-review}
The accompanying Wolfram script uses the exact rational form
$Q=\operatorname{diag}(3,5,7)$ and independent frame columns
$(1,1,0)^T,(0,1,1)^T$. It obtains
\[
 B=\begin{pmatrix}8&5\\5&12\end{pmatrix},\quad
 \det B=71,\quad
 Q(r,r)=\frac{105}{71}(a-b+c)^2 .
\]
It checks projector/convolution coefficients, residual orthogonality,
the Schur identity, symbolic alteration scaling, and all 18 sign
instances $1\le p\le6$, $2p+1\le M\le2p+3$.
The $M=5$, primitive K3-diagonal norm $21$ discriminator gives
doubled diagonal index $-21$, consistent with (11).
These are exact finite/symbolic consistency checks; the proofs
above establish the general identities.

Self-review covered nonprimitive alteration pullbacks, the two
independent duality signs, empty frames, dependence versus Gram
invertibility, rational versus integral perfect objects, and the
difference between a native kernel and a native target input.
No independent agent review or new Lean compilation was performed.
The new results are (5)--(13), the explicit actual return (16),
and the exact partial-trace period identity.
The universal native-source existence (17), and therefore the
original all-$X$ Hodge goal, remain unresolved here.
\end{document}
